\honest{Is Reality Ugly?}{Eliezer Yudkowsky}{2008}

Yesterday I showed that the cubes, 1, 8, 27, 64, 125, have first differences with no obvious pattern, but third differences that are all 6. That was a handpicked example. Mathematics starts from a few axioms and is closed; perhaps the real world is uglier. \nb{``Expecting Beauty,'' the post referred to, is not in this book.}

That question made sense two and a half thousand years ago. Heraclitus said all is fire, Thales water, Pythagoras number. ``Score: Heraclitus 0 Thales 0 Pythagoras 1.'' Beneath the surface there is a simple, exact level, which we call physics. Our best theories are mathematical and not yet unified, probably because they are surface forms of deeper mathematics, ``So by far the best guess is that reality is made of math.'' \nb{That physics is described by mathematics is what the paragraph argues. That the world is a mathematical structure is a stronger claim, Max Tegmark's hypothesis, and a controversial one. The post does not argue for it, and does not need it.}

Why, then, can we not predict tigers? I count three sources of uncertainty. First, we may not know the fundamental laws, but physicists need particle accelerators even to see that gap. Second, we may not have the computing power to work out what we know: we cannot predict a protein's shape from its amino acids. \nb{Since 2020, AlphaFold has, by learning from known structures rather than from physics.} Third, following a book by Nick Bostrom, we may not know where we are. A little person standing on one of the cubes, who knows all about the sequence, still has to look down to learn which cube is theirs. \nb{The same point is David Lewis's case of two gods (1979), who know every proposition true at their world but not which god they are.}

So real messiness comes from logical and indexical uncertainty. In biology we have already found the order, and it is too deep to help. ``Because'' of that, ``we can guess that we won't find any additional secret patterns'' that will make biology as easy as the cubes. \nb{The reasons given show only that biology cannot be derived from physics. Thermodynamics found simple laws for systems whose molecular detail no one can compute.} Likewise, no one will ever predict certain quantum results, ``only because our fundamental theory tells us quite definitely that different versions of us will see different results.'' \nb{That is the many-worlds reading, which, as the introduction to this book says, a number of physicists reject. On any standard reading the theory gives only probabilities, so the conclusion holds without it.}

The world is a ``perfectly regular, deterministic, and very large'' mathematical object. Uncertainty is in the map, not the territory. Even the blogosphere is as orderly as the cubes: ``the Internet is not a big muck... it's a series of cubes.''
